\documentclass[12pt]{article}

\pagestyle{empty}
\setlength{\topmargin}{0in}
\setlength{\headheight}{0in}
\setlength{\topsep}{0in}
\setlength{\textheight}{9in}
\setlength{\oddsidemargin}{0in}
\setlength{\evensidemargin}{0in}
\setlength{\textwidth}{6.5in}

\usepackage{palatino,graphics,amsmath,amssymb,enumitem}

\newcommand{\ds}{\displaystyle}
\newcommand{\vs}[1]{\vspace{#1in}}
\renewcommand{\vss}[1]{\vspace*{#1in}}
\newcommand{\bvec}{{\mathbf b}}
\newcommand{\cvec}{{\mathbf c}}
\newcommand{\dvec}{{\mathbf d}}
\newcommand{\evec}{{\mathbf e}}
\newcommand{\fvec}{{\mathbf f}}
\newcommand{\qvec}{{\mathbf q}}
\newcommand{\uvec}{{\mathbf u}}
\newcommand{\vvec}{{\mathbf v}}
\newcommand{\wvec}{{\mathbf w}}
\newcommand{\xvec}{{\mathbf x}}
\newcommand{\yvec}{{\mathbf y}}
\newcommand{\zvec}{{\mathbf y}}
\newcommand{\zerovec}{{\mathbf 0}}
\newcommand{\real}{{\mathbb R}}
\newcommand{\twovec}[2]{\left[\begin{array}{r}#1 \\ #2
    \end{array}\right]}
\newcommand{\ctwovec}[2]{\left[\begin{array}{c}#1 \\ #2
   \end{array}\right]}
\newcommand{\threevec}[3]{\left[\begin{array}{r}#1 \\ #2 \\ #3
  \end{array}\right]}
\newcommand{\cthreevec}[3]{\left[\begin{array}{c}#1 \\ #2 \\ #3
    \end{array}\right]}
\newcommand{\fourvec}[4]{\left[\begin{array}{r}#1 \\ #2 \\ #3 \\ #4
    \end{array}\right]}
\newcommand{\cfourvec}[4]{\left[\begin{array}{c}#1 \\ #2 \\ #3 \\ #4
    \end{array}\right]}
\newcommand{\mattwo}[4]{\left[\begin{array}{rr}#1 \amp #2 \\ #3 \amp #4 \\ \end{array}\right]}
\renewcommand{\span}[1]{\text{Span}\{#1\}}
\newcommand{\bcal}{{\cal B}}
\newcommand{\ccal}{{\cal C}}
\newcommand{\scal}{{\cal S}}
\newcommand{\wcal}{{\cal W}}
\newcommand{\ecal}{{\cal E}}
\newcommand{\coords}[2]{\left\{#1\right\}_{#2}}
\newcommand{\gray}[1]{\color{gray}{#1}}
\newcommand{\lgray}[1]{\color{lightgray}{#1}}
\newcommand{\rank}{\text{rank}}
\newcommand{\col}{\text{Col}}
\newcommand{\nul}{\text{Nul}}
\newcommand{\laspan}[1]{\text{Span}\{#1\}}

\begin{document}

\noindent
{\bf Mathematics 204} \\ 
{\bf Review}

\bigskip
\begin{enumerate}
\item Suppose that
  $$A = \begin{bmatrix}
    2 & -4 & -1 \\
    1 & -2 & 3 \\
    -1 & 2 & 4 \\
  \end{bmatrix},
  \hspace*{24pt}
  \bvec=\threevec077. \hspace*{24pt}
  $$
  
  Write the solution set of the equation $A\xvec = \bvec$ in vector
  parametric form.

  \vs{1.5}
  Remember that $\col(A)$ is the subspace that is the span of the
  columns of $A$.  Find a basis for $\col(A)$ and explain why your set
  is a basis.

  \vs{1}
  Is the vector $\bvec$ in $\col(A)$?  If so, write it as a linear
  combination of basis vectors.

  \vs{1.2}
  What is the dimension of $\col(A)$ and for which $p$ is
  $\col(A)$ a subspace of $\real^p$?

  \vs{1.2}
  Is the matrix $A$ invertible?  Explain your response.

  \vs{1}
  \newpage
  If a vector $\bvec$ is in $\col(A)$, what can you say about the set
  of solutions to the equation $A\xvec = \bvec$?

  \vs{1}
  
\item We earlier defined the {\em rank} of a matrix $A$ to be the
  number of pivot positions in $A$.  
  Give an example of a $3\times 4$ matrix for which $\rank(A) = 3$.
  What does this say about the columns of $A$? 

  \vs{1}
%  If $\rank(A) = 2$, what can you say about $\col(A)$?

  Give an example of a $3\times 4$ matrix whose entries are all
  nonzero for which $\rank(A)=1$.

  \vs{1}

  Suppose that $A$ is a $10\times15$ matrix.  What is the largest
  possible value for $\rank(A)$?

  \vs{1}
  If $\rank(A)$ is this largest possible value, what can you say about
  $\col(A)$ and what can you say about $\nul(A)$?

  \vs{1}
  If $B$ is a $12\times10$ matrix, what is the largest possible
  value of $\rank(B)$?

  \vs{1}
  \newpage
  If $\rank(B)$ is this largest possible value, what can you say about
  $\col(B)$ and what can you say about $\nul(B)$?

  % \vs{1}
  % If $C$ is a $14\times n$ matrix and $\col(C)=\real^{14}$, what can
  % you conclude about $n$?
  \vs{1}
\item Suppose that $A$ is a $3\times2$ matrix and that
  $A\twovec10 = \threevec{-2}{1}4$ and $A\twovec01 = \threevec11{-2}$.
  What is $A\twovec2{-3}$.

  \vs{1}
  Suppose that $B$ is a $2\times2$ matrix and that
  $$
  B\twovec23 = \twovec10,~~~B\twovec11 = \twovec01.
  $$
  What is $B^{-1}$?

  \vs{1}
  Find a solution to the equation $B\xvec=\twovec4{-1}$.

  \vs{1}
  What is the matrix $B$?

  \vs{1}
  \newpage
\item For what values of $k$ is the matrix
  $\begin{bmatrix}
    3 & 2 & 0 \\
    k & -1 & 1 \\
    1 & 2 & 2 \\
  \end{bmatrix}$ invertible?
  
  \vs{1.5}
\item What does it mean for $\vvec$ to be an eigenvector of a matrix
  $A$?

  \vs{1}

  Suppose that $A$ is a $2\times2$ matrix having eigenvectors
  $\vvec_1 = \twovec21$ and $\vvec_2=\twovec31$ and eigenvalues
  $\lambda_1=4$ and $\lambda_2=-2$.

  Write the vector $\twovec10$ as a linear combination of
  eigenvectors.

  \vs{1}
  Find the vector $A\twovec10$.

  \vs{1}
  In the same way, write the vector $\twovec01$ as a linear
  combination of eigenvectors.  Then find the vector $A\twovec01$.

  \vs{1.5}
  What is the matrix $A$?

  \vs{1}
  Is $A$ diagonalizable?  If so, find the matrices $P$ and $D$ and
  verify that $A=PDP^{-1}$.
    
  
  

\end{enumerate}
\end{document}



  



  

